% Options for packages loaded elsewhere
% Options for packages loaded elsewhere
\PassOptionsToPackage{unicode}{hyperref}
\PassOptionsToPackage{hyphens}{url}
\PassOptionsToPackage{dvipsnames,svgnames,x11names}{xcolor}
%
\documentclass[
]{agujournal2019}
\usepackage{xcolor}
\usepackage{amsmath,amssymb}
\setcounter{secnumdepth}{5}
\usepackage{iftex}
\ifPDFTeX
  \usepackage[T1]{fontenc}
  \usepackage[utf8]{inputenc}
  \usepackage{textcomp} % provide euro and other symbols
\else % if luatex or xetex
  \usepackage{unicode-math} % this also loads fontspec
  \defaultfontfeatures{Scale=MatchLowercase}
  \defaultfontfeatures[\rmfamily]{Ligatures=TeX,Scale=1}
\fi
\usepackage{lmodern}
\ifPDFTeX\else
  % xetex/luatex font selection
\fi
% Use upquote if available, for straight quotes in verbatim environments
\IfFileExists{upquote.sty}{\usepackage{upquote}}{}
\IfFileExists{microtype.sty}{% use microtype if available
  \usepackage[]{microtype}
  \UseMicrotypeSet[protrusion]{basicmath} % disable protrusion for tt fonts
}{}
\makeatletter
\@ifundefined{KOMAClassName}{% if non-KOMA class
  \IfFileExists{parskip.sty}{%
    \usepackage{parskip}
  }{% else
    \setlength{\parindent}{0pt}
    \setlength{\parskip}{6pt plus 2pt minus 1pt}}
}{% if KOMA class
  \KOMAoptions{parskip=half}}
\makeatother
% Make \paragraph and \subparagraph free-standing
\makeatletter
\ifx\paragraph\undefined\else
  \let\oldparagraph\paragraph
  \renewcommand{\paragraph}{
    \@ifstar
      \xxxParagraphStar
      \xxxParagraphNoStar
  }
  \newcommand{\xxxParagraphStar}[1]{\oldparagraph*{#1}\mbox{}}
  \newcommand{\xxxParagraphNoStar}[1]{\oldparagraph{#1}\mbox{}}
\fi
\ifx\subparagraph\undefined\else
  \let\oldsubparagraph\subparagraph
  \renewcommand{\subparagraph}{
    \@ifstar
      \xxxSubParagraphStar
      \xxxSubParagraphNoStar
  }
  \newcommand{\xxxSubParagraphStar}[1]{\oldsubparagraph*{#1}\mbox{}}
  \newcommand{\xxxSubParagraphNoStar}[1]{\oldsubparagraph{#1}\mbox{}}
\fi
\makeatother


\usepackage{longtable,booktabs,array}
\newcounter{none} % for unnumbered tables
\usepackage{calc} % for calculating minipage widths
% Correct order of tables after \paragraph or \subparagraph
\usepackage{etoolbox}
\makeatletter
\patchcmd\longtable{\par}{\if@noskipsec\mbox{}\fi\par}{}{}
\makeatother
% Allow footnotes in longtable head/foot
\IfFileExists{footnotehyper.sty}{\usepackage{footnotehyper}}{\usepackage{footnote}}
\makesavenoteenv{longtable}
\usepackage{graphicx}
\makeatletter
\newsavebox\pandoc@box
\newcommand*\pandocbounded[1]{% scales image to fit in text height/width
  \sbox\pandoc@box{#1}%
  \Gscale@div\@tempa{\textheight}{\dimexpr\ht\pandoc@box+\dp\pandoc@box\relax}%
  \Gscale@div\@tempb{\linewidth}{\wd\pandoc@box}%
  \ifdim\@tempb\p@<\@tempa\p@\let\@tempa\@tempb\fi% select the smaller of both
  \ifdim\@tempa\p@<\p@\scalebox{\@tempa}{\usebox\pandoc@box}%
  \else\usebox{\pandoc@box}%
  \fi%
}
% Set default figure placement to htbp
\def\fps@figure{htbp}
\makeatother


% definitions for citeproc citations
\NewDocumentCommand\citeproctext{}{}
\NewDocumentCommand\citeproc{mm}{%
  \begingroup\def\citeproctext{#2}\cite{#1}\endgroup}
\makeatletter
 % allow citations to break across lines
 \let\@cite@ofmt\@firstofone
 % avoid brackets around text for \cite:
 \def\@biblabel#1{}
 \def\@cite#1#2{{#1\if@tempswa , #2\fi}}
\makeatother
\newlength{\cslhangindent}
\setlength{\cslhangindent}{1.5em}
\newlength{\csllabelwidth}
\setlength{\csllabelwidth}{3em}
\newenvironment{CSLReferences}[2] % #1 hanging-indent, #2 entry-spacing
 {\begin{list}{}{%
  \setlength{\itemindent}{0pt}
  \setlength{\leftmargin}{0pt}
  \setlength{\parsep}{0pt}
  % turn on hanging indent if param 1 is 1
  \ifodd #1
   \setlength{\leftmargin}{\cslhangindent}
   \setlength{\itemindent}{-1\cslhangindent}
  \fi
  % set entry spacing
  \setlength{\itemsep}{#2\baselineskip}}}
 {\end{list}}
\usepackage{calc}
\newcommand{\CSLBlock}[1]{\hfill\break\parbox[t]{\linewidth}{\strut\ignorespaces#1\strut}}
\newcommand{\CSLLeftMargin}[1]{\parbox[t]{\csllabelwidth}{\strut#1\strut}}
\newcommand{\CSLRightInline}[1]{\parbox[t]{\linewidth - \csllabelwidth}{\strut#1\strut}}
\newcommand{\CSLIndent}[1]{\hspace{\cslhangindent}#1}



\setlength{\emergencystretch}{3em} % prevent overfull lines

\providecommand{\tightlist}{%
  \setlength{\itemsep}{0pt}\setlength{\parskip}{0pt}}



 


\usepackage{url} %this package should fix any errors with URLs in refs.
\usepackage{lineno}
\usepackage[inline]{trackchanges} %for better track changes. finalnew option will compile document with changes incorporated.
\usepackage{soul}
\linenumbers
\makeatletter
\@ifpackageloaded{caption}{}{\usepackage{caption}}
\AtBeginDocument{%
\ifdefined\contentsname
  \renewcommand*\contentsname{Table of contents}
\else
  \newcommand\contentsname{Table of contents}
\fi
\ifdefined\listfigurename
  \renewcommand*\listfigurename{List of Figures}
\else
  \newcommand\listfigurename{List of Figures}
\fi
\ifdefined\listtablename
  \renewcommand*\listtablename{List of Tables}
\else
  \newcommand\listtablename{List of Tables}
\fi
\ifdefined\figurename
  \renewcommand*\figurename{Figure}
\else
  \newcommand\figurename{Figure}
\fi
\ifdefined\tablename
  \renewcommand*\tablename{Table}
\else
  \newcommand\tablename{Table}
\fi
}
\@ifpackageloaded{float}{}{\usepackage{float}}
\floatstyle{ruled}
\@ifundefined{c@chapter}{\newfloat{codelisting}{h}{lop}}{\newfloat{codelisting}{h}{lop}[chapter]}
\floatname{codelisting}{Listing}
\newcommand*\listoflistings{\listof{codelisting}{List of Listings}}
\makeatother
\makeatletter
\makeatother
\makeatletter
\@ifpackageloaded{caption}{}{\usepackage{caption}}
\@ifpackageloaded{subcaption}{}{\usepackage{subcaption}}
\makeatother
\usepackage{bookmark}
\IfFileExists{xurl.sty}{\usepackage{xurl}}{} % add URL line breaks if available
\urlstyle{same}
% fallback for those not using the hyperref driver hyperxmp:
\makeatletter
\@ifundefined{xmpquote}{\newcommand{\xmpquote}[1]{#1}}{}
\makeatother
\hypersetup{
  pdftitle={Assessing the influence of variable annual sampling effort on official national forest inventory statistics},
  pdfauthor={George C. Gaines III; Bryce Frank; KaDonna Randolph; Jim Westfall},
  pdfkeywords={\xmpquote{Forest Inventory and Analysis}},
  colorlinks=true,
  linkcolor={blue},
  filecolor={Maroon},
  citecolor={Blue},
  urlcolor={Blue},
  pdfcreator={LaTeX via pandoc}}


\journalname{Forest Ecology and Management}

\draftfalse

\begin{document}
\title{Assessing the influence of variable annual sampling effort on
official national forest inventory statistics}

\authors{George C. Gaines III\affil{1}, Bryce Frank\affil{2}, KaDonna
Randolph\affil{3}, Jim Westfall\affil{4}}
\affiliation{1}{USDA Forest Service, Rocky Mountain Research
Station, Missoula, MT, }\affiliation{2}{USDA Forest Service, Pacific
Northwest Research Station, XXX, WA, }\affiliation{3}{USDA Forest
Service, Southern Research
Station,, Knoxville, TN, }\affiliation{4}{USDA Forest Service, Northern
Research Station, XXX, PA, }
\correspondingauthor{George C. Gaines III}{george.gaines@usda.gov}


\begin{abstract}
In this study we\ldots{}
\end{abstract}





\section{Introduction}\label{introduction}

The USFS Forest Inventory and Analysis (FIA) Program surveys and
analyzes the forests of the United States through a nationwide network
of permanent sample locations. These plot locations are distributed
across all lands with a spatial sampling intensity of about 1 location
per 2400 ha according to an unaligned, quasi-systematic equal
probability sample design. This design forces desirable spatial
dispersion while limiting the potential for alignment of sample
locations with periodic population features. Statistical inference on
forest population status and change parameters is enabled by
measurements of forest vegetation traits made at these locations by
field crews.

Historically, FIA sample locations in each state were visited
periodically as time, budgets, and workforce capacity allowed. This
approach offered logistical flexibility but complicated analysis and
reporting timelines. In the early 2000's the program shifted from
periodic to annual inventory protocols calling for approximately equal
numbers of plot remeasurements each year on schedules intended to be
fixed and consistent within states. Under the annual inventory, each
plot location in a state is assigned to one of a set of panels, and all
plots assigned to a panel are designated for measurement in the same
year. Once the final panel in an inventory cycle is complete,
remeasurements at sample locations in the first panel begin the
following year. For more information regarding the FIA sample design see
(Westfall, Coulston, et al., 2022, Chapter 2).

In general, FIA requires that a panel be complete before the associated
raw tabular data are pushed to its public-facing databases, and requires
a complete inventory cycle within a state to publish official
state-scale estimates of forest parameters in fulfillment of
Congressional reporting obligations. Unfortunately, in practice,
adherence to the annual design described above has proven very
challenging or impossible for the FIA Program. Instead, delayed and
skipped field measurements cause unintentional lengthening of inventory
cycles and/or survey nonresponse (see, e.g., Särndal \& Lundström,
2005), at rates that vary regionally and over time. The underlying
causes of delayed and skipped field measurements are varied, and can be
chronic (e.g., denied access to private lands (Westfall, Schroeder, et
al., 2022), logistical or administrative challenges in remote regions
(Affleck et al., 2026), persistent workforce and budgetary constraints
(Conroy, 2026)) or result from discrete events (e.g., the COVID19
pandemic). Better understanding of patterns in late and skipped
measurements (Conroy, 2026), and the investigation of new sampling and
estimation techniques to ameliorate the statistical consequences
nonresponse (Westfall, 2022; Westfall \& Patterson, 2026), are active
research topics.

Of immediate practical concern for FIA and its external partners are the
strain on field staff and delays to tabular data publication that result
from delayed field measurements. Methods have been proposed to control
the measurement delay process by disaggregating FIA's annual panel
structure while still preserving the spatial dispersion of the FIA
sample network within field season years. Panel disaggregation enables
intentional alignment of workloads with anticipated data acquisition
workforce capacity in each field season year through the application of
variable annual field sampling effort within individual states.

Though the logistical advantages of variable annual effort (VAE) are
profound (i.e., it can help ensure the long-term viability of the
inventory in the context of funding and staffing volatility), the
statistical consequences of the redistribution of partial panels within
inventory cycles, and of extended and variable inventory cycle lengths,
are unknown across states and epochs. In particular, a need exists to
understand and communicate to stakeholders how adoption of VAE protocols
may influence the statistical properties of forest parameter estimates
published by FIA in fulfillment of Congressional reporting imperatives.

In this study, we seek to\ldots{}

\section{Data \& Methods}\label{sec-data-methods}

\begin{figure}[H]

\centering{

\pandocbounded{\includegraphics[keepaspectratio,alt={A scatterplot of earthquake locations plotting latitude against longitude.}]{index_vae_files/figure-latex/notebooks-explore-earthquakes-fig-spatial-plot-output-1.png}}

}

\caption{\label{fig-spatial-plot}Locations of earthquakes on La Palma
since 2017}

\end{figure}%

\textsubscript{Source:
\href{https://gcgainescubed.github.io/manuscript_template/notebooks/explore-earthquakes-preview.html\#cell-fig-spatial-plot}{Explore
Earthquakes}}

\section{Conclusion}\label{conclusion}

\section*{References}\label{references}
\addcontentsline{toc}{section}{References}

\protect\phantomsection\label{refs}
\begin{CSLReferences}{1}{0}
\vspace{1em}

\bibitem[\citeproctext]{ref-affleck2026inventory}
Affleck, D. L., Gaines III, G. C., Conroy, C., Santino, G., \& Larson,
A. J. (2026). Inventory measurement rates lag in wilderness areas across
the western USA. \emph{Forest Science}, \emph{72}(1), 95--102.

\bibitem[\citeproctext]{ref-conroy_thesis}
Conroy, C. (2026). Elevated FIA nonresponse rates in wilderness:
Investigating barriers and potential solutions.

\bibitem[\citeproctext]{ref-sarndal_nonresponse}
Särndal, C.-E., \& Lundström, S. (2005). \emph{Estimation in surveys
with nonresponse}. John Wiley \& Sons.

\bibitem[\citeproctext]{ref-westfall2022estimation}
Westfall, J. A. (2022). An estimation method to reduce complete and
partial nonresponse bias in forest inventory. \emph{European Journal of
Forest Research}, \emph{141}(5), 901--907.

\bibitem[\citeproctext]{ref-westfall2026double}
Westfall, J. A., \& Patterson, P. L. (2026). Double ratio estimation
within a design-based nonresponse bias mitigation strategy.
\emph{Forestry: An International Journal of Forest Research},
\emph{99}(1), cpaf032.

\bibitem[\citeproctext]{ref-westfall2022spatial}
Westfall, J. A., Schroeder, T. A., McCollum, J. M., \& Patterson, P. L.
(2022). A spatial and temporal assessment of nonresponse in the national
forest inventory of the US. \emph{Environmental Monitoring and
Assessment}, \emph{194}(8), 530.

\bibitem[\citeproctext]{ref-gb2}
Westfall, J. A., Coulston, J. W., Moisen, G. G., \& Andersen, H.-E.
(2022). Sampling and estimation documentation for the enhanced forest
inventory and analysis program: 2022. \emph{Gen. Tech. Rep. NRS-GTR-207.
Madison, WI: US Department of Agriculture, Forest Service, Northern
Research Station. 129 p.}, \emph{207}, 1--129.

\end{CSLReferences}




\end{document}
